% Source: https://artofproblemsolving.com/community/c3h1871697_function
% Stable identity: aops-topic-1871697-q01
$h^{-1}(4)$ is the input $w$ with $h(w)=4$. Write $w=2z+3$, so $z=\dfrac{w-3}{2}$ and
\[
h(w)=\dfrac{5-(w-3)}{4+\dfrac{w-3}{2}}=\dfrac{2(8-w)}{w+5}\qquad(w\ne-5).
\]
Then $h(w)=4$ means $2(8-w)=4(w+5)$, i.e.\ $16-2w=4w+20$, so $w=-\dfrac{2}{3}$. This is the only solution, so $h$ takes the value $4$ exactly once and
\[
h^{-1}(4)=\boxed{-\dfrac{2}{3}}.
\]
(Equivalently: $\dfrac{5-2z}{4+z}=4$ gives $z=-\dfrac{11}{6}$, and then $2z+3=-\dfrac{11}{3}+3=-\dfrac{2}{3}$. Check: $h\left(-\dfrac{2}{3}\right)=\dfrac{2\cdot\dfrac{26}{3}}{\dfrac{13}{3}}=4$.)
